\chapter{Mirtazapine (active ingredient)}
\index{Mirtazapine (active ingredient)}

\section*{Definition and Overview}
Mirtazapine is a noradrenergic and specific serotonergic antidepressant (NaSSA) utilised primarily for the treatment of Major Depressive Disorder (MDD) in adults. Mechanistically, it differs from common antidepressants like Selective Serotonin Reuptake Inhibitors (SSRIs) and Serotonin-Noradrenaline Reuptake Inhibitors (SNRIs) by acting as an antagonist at central presynaptic $\alpha_2$-adrenergic receptors and specific postsynaptic serotonin receptors, rather than inhibiting monoamine reuptake. Additionally, mirtazapine exhibits potent antagonism at histamine H$_1$ receptors, which account for its prominent sedative and appetite-stimulating properties. It is frequently prescribed when depression is accompanied by severe insomnia, anxiety, or significant weight loss, and it has a lower incidence of sexual dysfunction and gastrointestinal side effects compared to SSRIs.

\section*{Epidemiology}
Mirtazapine is one of the most commonly prescribed antidepressants globally. In many countries, it is widely dispensed under the public medicines subsidy (public medicines subsidy) for the management of moderate-to-severe depression. Due to its sedative and orexigenic (appetite-stimulating) effects, it is particularly favoured in the elderly population (geriatric depression) to address age-associated insomnia and anorexia (failure to thrive). While approved primarily for MDD, it is also frequently used off-label for generalised anxiety disorder, chronic insomnia, post-traumatic stress disorder, and tension-type headaches.

\section*{Aetiology and Pathophysiology}
The therapeutic and adverse effects of mirtazapine are directly related to its unique receptor binding profile in the central and peripheral nervous systems.
\begin{itemize}
    \item \textbf{Presynaptic $\alpha_2$-adrenergic receptor antagonism:} Mirtazapine blocks central presynaptic $\alpha_2$ auto-receptors on noradrenergic neurones and $\alpha_2$ hetero-receptors on serotonergic neurones. This blockade removes the negative feedback mechanism, leading to increased release of both noradrenaline and serotonin (5-HT) into the synaptic cleft.
    \item \textbf{Postsynaptic 5-HT$_2$ and 5-HT$_3$ receptor antagonism:} By blocking these specific serotonin receptors, mirtazapine redirects serotonergic neurotransmission exclusively to 5-HT$_{1A}$ receptors. This targeted action avoids the side effects commonly associated with non-specific 5-HT$_2$ and 5-HT$_3$ stimulation, such as anxiety, insomnia, sexual dysfunction, nausea, and headache.
    \item \textbf{Histamine H$_1$ receptor antagonism:} Mirtazapine has a very high affinity for H$_1$ receptors, causing potent sedation, drowsiness, and weight gain.
    \item \textbf{Inverse relationship between dose and sedation:} Paradoxically, mirtazapine is often more sedating at lower doses (e.g. 7.5 mg to 15 mg) than at higher doses (30 mg to 45 mg). At higher doses, the increased noradrenergic neurotransmission counteracts the histaminergic sedative effect.
    \item \textbf{Minimal anticholinergic activity:} Mirtazapine has very low affinity for muscarinic cholinergic and $\alpha_1$-adrenergic receptors, resulting in fewer side effects like dry mouth, urinary retention, and orthostatic hypotension compared to tricyclic antidepressants (TCAs).
\end{itemize}

\section*{Clinical Features}
The clinical features of mirtazapine therapy encompass both the resolution of depressive symptoms and the presentation of its characteristic side effect profile.
\begin{itemize}
    \item \textbf{Therapeutic response:} Gradual improvement in mood, reduction in anxiety, and a rapid, noticeable improvement in sleep quality and sleep architecture, often occurring within the first week of starting therapy.
    \item \textbf{Somnolence and fatigue:} Significant daytime drowsiness, lethargy, and prolonged sleep, which are most pronounced during the first few weeks of treatment.
    \item \textbf{Increased appetite and weight gain:} A marked increase in food cravings, particularly for carbohydrates, leading to rapid weight gain.
    \item \textbf{Antihistaminic and mild anticholinergic effects:} Dry mouth, constipation, increased thirst, and occasional lightheadedness upon standing (dizziness).
    \item \textbf{Blood disorders (Neutropenia):} A rare but serious clinical feature characterised by a drop in white blood cell count, presenting as sudden fever, sore throat, or mouth ulcers.
    \item \textbf{Serotonin syndrome features:} Agitation, confusion, rapid heart rate, dilated pupils, muscle twitching, shivering, sweating, and hyperreflexia.
\end{itemize}

\section*{Differential Diagnosis}
Clinicians must differentiate mirtazapine's side effects and complications from the underlying psychiatric illness and other medical conditions.
\begin{itemize}
    \item \textbf{Mirtazapine-induced somnolence vs.\ depressive apathy or hypersomnia:} Distinguishing drug-induced daytime sleepiness from a vegetative symptom of severe depression, which would require a dose increase rather than decrease.
    \item \textbf{Drug-induced weight gain vs.\ metabolic or endocrine disorders:} Differentiating mirtazapine-associated weight gain from hypothyroidism, Cushing's syndrome, or atypical depression.
    \item \textbf{Serotonin syndrome vs.\ Neuroleptic Malignant Syndrome (NMS) vs.\ Anticholinergic toxicity:} Distinguishing between drug-induced toxicities. Serotonin syndrome is characterized by hyperreflexia and clonus, whereas NMS features "lead-pipe" muscle rigidity, and anticholinergic toxicity is marked by dry skin and absent bowel sounds.
    \item \textbf{Drug-induced neutropenia vs.\ acute viral infection:} Differentiating a mild viral illness from life-threatening agranulocytosis in a patient on mirtazapine who presents with a sore throat and fever.
\end{itemize}

\section*{When to Seek Medical Care}
Patients taking mirtazapine should consult their general practitioner or psychiatrist if they notice worsening depression, suicidal thoughts, unusual changes in behaviour, or signs of an infection such as a sudden fever, sore throat, or mouth ulcers.

\textbf{Contact your local emergency services for:}
\begin{itemize}
    \item Symptoms of serotonin syndrome, including high fever, severe sweating, rapid heart rate, muscle rigidity, shivering, severe confusion, and uncontrollable muscle spasms.
    \item Acute overdose, presenting with severe drowsiness, disorientation, rapid heart rate, or seizures.
    \item Severe chest pain, shortness of breath, or palpitations (suspected QTc prolongation or arrhythmia).
    \item Signs of a severe allergic reaction (anaphylaxis), such as swelling of the face, lips, or throat, or a spreading, blistering skin rash (Stevens-Johnson syndrome).
\end{itemize}

\section*{Diagnosis and Investigations}
Monitoring patients on mirtazapine involves psychological assessment and targetted laboratory testing.
\begin{itemize}
    \item \textbf{Full Blood Count (FBC):} Indicated immediately if the patient develops a fever, sore throat, stomatitis, or other signs of infection, to check for neutropenia or agranulocytosis (which typically occurs within 4 to 6 weeks of initiation).
    \item \textbf{Metabolic panel (Lipids and Glucose):} Periodic monitoring of fasting blood glucose and lipid profile (cholesterol and triglycerides), as mirtazapine can induce hyperlipidaemia and weight gain.
    \item \textbf{Electrocardiogram (ECG):} Performed baseline and periodically in patients with risk factors for QTc prolongation (e.g. cardiac disease, co-administration of other QTc-prolonging drugs).
    \item \textbf{Clinical psychiatric monitoring:} Structured assessment of suicide risk, depression severity, and adherence using validated tools (e.g. PHQ-9, C-SSRS).
\end{itemize}

\section*{Management}
Management of mirtazapine therapy involves careful dosing, side effect management, and structured discontinuation protocols.
\begin{itemize}
    \item \textbf{Dosing and administration:} Typically initiated at 15 mg or 30 mg orally once daily, taken strictly at bedtime due to its sedative effects. The dose may be titrated up to a maximum of 45 mg daily.
    \item \textbf{Managing sedation:} Educating patients that lower doses can cause more sedation, and that this effect often improves with dose escalation to 30 mg or 45 mg, or with continued use over 2 to 4 weeks.
    \item \textbf{Gradual discontinuation (tapering):} Avoiding sudden cessation, which can trigger antidepressant discontinuation syndrome (characterised by anxiety, insomnia, headache, nausea, and dizziness). The dose should be tapered gradually over 2 to 4 weeks.
    \item \textbf{Managing agranulocytosis:} If neutropenia (neutrophils $<1.0 \times 10^9$/L) is confirmed, mirtazapine must be discontinued immediately and permanently, and the patient monitored closely.
    \item \textbf{Serotonin syndrome protocol:} Discontinuing mirtazapine, administering intravenous benzodiazepines for sedation and muscle relaxation, supportive cooling, and administering serotonin antagonists (e.g. cyproheptadine).
\end{itemize}

\section*{Complications}
The primary complications of mirtazapine relate to metabolic changes, bone marrow suppression, and drug interactions.
\begin{itemize}
    \item \textbf{Metabolic syndrome:} Significant weight gain, hyperlipidaemia, and insulin resistance increasing the risk of type 2 diabetes and cardiovascular disease.
    \item \textbf{Agranulocytosis and neutropenia:} Rare but potentially fatal bone marrow suppression, rendering the patient highly vulnerable to severe systemic infections.
    \item \textbf{Serotonin syndrome:} A life-threatening toxic state occurring primarily when mirtazapine is co-administered with other serotonergic agents (e.g. MAOIs, SSRIs, tramadol).
    \item \textbf{QTc interval prolongation:} Increased risk of torsades de pointes and ventricular fibrillation, especially in overdose or when combined with other QTc-prolonging drugs.
\end{itemize}

\section*{Prognosis and Outcomes}
The prognosis for patients treated with mirtazapine is highly favourable, with response and remission rates for depression comparable to or slightly higher than those of SSRIs. Sleep disturbances and anxiety often resolve rapidly within the first few days of therapy, which can improve treatment compliance. The majority of side effects, such as dry mouth and transient daytime somnolence, tend to subside within the first month. Weight gain can be long-lasting and require active lifestyle management.

\section*{Prevention and Self-Care}
Patients can manage the side effects of mirtazapine and ensure safe therapy through proactive self-care habits.
\begin{itemize}
    \item \textbf{Strict bedtime dosing:} Taking mirtazapine immediately before going to bed to utilise its sedative properties to aid sleep and minimise daytime somnolence.
    \item \textbf{Avoid driving when drowsy:} Refraining from driving, operating heavy machinery, or performing dangerous tasks if feeling drowsy or dizzy.
    \item \textbf{Monitor food intake and weight:} Adopting a balanced, calorie-controlled diet and engaging in regular physical exercise to manage the drug-induced increase in appetite and prevent excessive weight gain.
    \item \textbf{Avoid alcohol:} Refraining from drinking alcohol while taking mirtazapine, as it significantly enhances the sedative effects and increases the risk of accidents.
    \item \textbf{Report signs of infection:} Reporting any fever, sore throat, or mouth ulcers to a doctor immediately, and request a blood test.
    \item \textbf{Maintain dental hygiene:} Using sugar-free lozenges, chewing gum, or saliva substitutes to manage dry mouth and prevent dental cavities.
\end{itemize}

\section*{Sources and Further Reading}
The following organisations provide authoritative information on mirtazapine and antidepressant management:
\begin{itemize}
    \item Australian Medicines Handbook (AMH). \textit{Mirtazapine monograph and clinical guidelines.} https://amh.org.au/
    \item NPS MedicineWise. \textit{Mirtazapine: Active ingredient information and resources.} https://www.nps.org.au/
    \item Royal Australian and New Zealand College of Psychiatrists (RANZCP). \textit{Clinical practice guidelines for the treatment of depression.} https://www.ranzcp.org/
    \item Mayo Clinic. \textit{Mirtazapine (Oral Route): Side effects and drug interactions.} https://www.mayoclinic.org/
    \item NHS. \textit{Mirtazapine: Antidepressant medicine overview.} https://www.nhs.uk/
\end{itemize}
